\documentclass[
oneside,							% einseitiges Layout
captions=signature,	  % caps immer unter Element (Bild, Grafik...)
openany,							% Kapitel können auch auf linken Seiten beginnen
listof=totoc				% Abb.- und Tab.verzeichnis im Inhalt	
]
{scrbook}


\usepackage[utf8]{inputenc}
\usepackage[T1]{fontenc}

\usepackage{mathptmx}
\usepackage[scaled=.90]{helvet}
\usepackage{courier}
					
\usepackage[ngerman]{babel}


\usepackage{qtree} % Baumstrukturen zeichnen

\usepackage{tabularx}
\usepackage{tabulary}
\usepackage{multirow}
\usepackage{colortbl}
\usepackage{booktabs}
\usepackage{listings}

%% Packages für Grafiken & Abbildungen %%%%%%%%%%%%%%%%%%%%%%

\usepackage{graphicx} 	% Zum Laden von Grafiken
\usepackage{xcolor}
\usepackage{color}
\usepackage{microtype} 	% visual enhancement randausgleich
\usepackage{float} 			% ermöglicht den H param bei Bilder 
												% für Plazierungen an der Position

\usepackage{rotating} % Text drehen

\usepackage[nogroupskip, nonumberlist, toc]{glossaries}
\usepackage{glossary-inline}
\usepackage[printonlyused]{acronym}

\usepackage{pifont} 	%symbol fonts like Windings
\usepackage{amsmath}	% Math font optimisation
\usepackage{amssymb} 	% Math symbols

\usepackage{tocloft} % Table of content - list of figures, tables - management


%% Bibliographiestil %%%%%%%%%%%%%%%%%%%%%%%%%%%%%%%%%%%%%%%%%%%%%%%%%%

\usepackage{natbib}
\renewcommand{\familydefault}{\sfdefault}


\setlength{\parindent}{0cm}			% Absatzeinzug
\setlength{\parskip}{6pt}				% Leerzeile zwischen Absätzen im Inhaltsverzeichnis
\textwidth=15.7cm								% Textbreite
\leftmargin=-3cm

\textheight=22.6cm							% Texthöhe
\topmargin=-0.7cm								% Rand (Vorsicht minus)
\headheight=0.2cm								% Höhe der Kopfzeile
\headsep=1.5cm									% Abstand von der Kopfzeile zum Text

%\setlength{\bibsep}{0.4\parskip}   	 	%Bibliographieabstand fest
\setlength{\bibsep}{0.3em plus 0.5em}	
%Bibliographieabstand variabel um Seitenumbruch zu unterdrücken

\setlength{\bibhang}{1cm}    			%Bibliographieeinzug
%\newcommand{\bibfont}{\normalsize}   	%Bibliographiezeichengröße
\renewenvironment{quote}{\list{}{\rightmargin0pt\leftmargin1cm}\item\relax}{\endlist} 
% neue \quote-Definition als globale Variabel für rechtsbündige Zitate

\setcitestyle{notesep={: }}
\setcitestyle{aysep={}}

%%%%%%%%%%%%%%%%%%%%%%%%%%%%%%%\setlist{\setlength{\itemsep}{0.8em plus 2em}, parsep=.65em, partopsep=1em} %listenformatierung

%%%%%% TOC Tiefe %%%%%%%%%%%%%%%%%%%%%%%%%%%%%%%%%%%%%
\setcounter{tocdepth}{3}

%%%%%%%% Nummerierungstiefe absätze %%%%%%%%%%%%%%%%%%
\setcounter{secnumdepth}{3}

%\setkomafont{section}{\LARGE}
%\setkomafont{subsection}{\Large}
%\setkomafont{subsubsection}{\large}
%\setkomafont{paragraph}{\large}
%\setkomafont{subparagraph}{\large}
% Die Überschriften sind in der Lokalen *.cls 
% Classendefinition definiert %%%%%


%% Abk. und Kapitälchen%%%%%%%%%%%%%%%%%%%%%%%%%%%%%%%%%%%%%%%%%%%%%%%%


\setglossarystyle{index}
\makenoidxglossaries
\setacronymstyle{sc-short-long}

\newacronym{fmx}{Formex}{Formalized Exchange of Electronic Publications}
\newacronym{cl}{cl}{Computerlinguistik der Universität Trier}
\newacronym{darpa}{darpa}{Defence Advanced Research Projects Agency}
\newacronym{bleu}{bleu}{Bilingual Evaluation Understudy}
\newacronym{nist}{nist}{National Institute of Standards and Technology}
\newacronym{hter}{hter}{Human-targeted Translation Edit Rate}
\newacronym{rbmt}{rbmt}{Rule Based Machine Translation -- Regelbasierte Maschinelle Übersetzung}
\newacronym{smt}{smt}{Statistical Machine Translation -- Statistische Maschinelle Übersetzung}
\newacronym{meteor}{meteor}{Metric for Evaluation of Translation with Explicit ORdering}
\newacronym{lepor}{lepor}{Length Penalty, Precision, n-gram Position difference Penalty and Recall}
\newacronym{amber}{amber}{A Modified Bleu Enhanced Ranking}

\newglossaryentry{glos:precision}{name=precision, description={Genauigkeit}}
\newglossaryentry{glos:recall}{name=precision, description={Trefferquote}}


%%% Trennungsregeln für das gesamte Dokument.

\hyphenation{Über-set-zungs-epo-che}
\hyphenation{Über-set-zungs-ver-fah-ren}
\hyphenation{Über-set-zung}
\hyphenation{Kor-re-la-tions-koef-fi-zient}
\hyphenation{Kor-re-la-tions-koef-fi-zien-ten}
\hyphenation{Rang-Kor-re-la-tions-koef-fi-zien-ten}


%%definition von quotation

\let\oldquotation\quotation
\renewcommand\quotation{\oldquotation\noindent\renewcommand{\baselinestretch}{1}\small\ignorespaces} 

%% Absolute Value definition %

\delimitershortfall-1sp
\newcommand\abs[1]{\left|#1\right|}
%%%%%%%%%%%%%%%%%%%%%%%%%%%%%%%%%%%%%%%%%%%%%%%

\providecommand\phantomsection{}

% Farbige unterteilung

\newcommand{\status}[2]{%
  \colorlet{author#1}{#2}% Create an author colour
  \expandafter\def\csname authoredby#1\endcsname{% Create author colour settings
    \renewcommand{\cftchapfont}{\bfseries\color{author#1}}% Chapter colour
    \renewcommand{\cftsecfont}{\color{author#1}}% Section colour
    \renewcommand{\cftsubsecfont}{\color{author#1}}% Subsection colour
	\renewcommand{\cftsubsubsecfont}{\color{author#1}}}% Subsection colour
	}

\makeatletter
\newcommand{\setstatus}[1]{\addtocontents{toc}{\protect\@nameuse{authoredby#1}}}%
\makeatother

\status{R}{red}% Author A will be coloured red
\status{Y}{orange}% Author B will be coloured blue
\status{G}{green}% Author B will be coloured blue


\begin{document}




%%%%%% T I T E L S E I T E %%%%%%%

\thispagestyle{empty}

\vspace*{-2cm} \hspace{-1cm}
{\includegraphics[height=2.5cm]{BILDER/logo.png}} \\ \vspace{3cm}
	\begin{center}
		\textsf{ \Huge{\textbf{Hausarbeit}}}\\ \vspace{1.5cm}
		\textsf{ \Huge{zu dem Thema}} \vspace{1.5cm}
			\begin{center}
				\textsf{ \Huge{\textbf{Vorlage zur Erstellung einer Hausarbeit im Fach CL}}}
			\end{center}
			\vfill \vfill \textsf{\Large \today}
	\end{center}
\vfill \hspace{-2cm}
\textsf{{\large
\begin{tabular}{ll}
		Vorname Name\\
		Email: mein.name@uni-trier.de\\
\end{tabular}
}}
\newpage




%% Erzeugung von Verzeichnissen %%%%%%%%%%%%%%%%%%%%%%%%%%%%%%%%%%%%%%%
\addtocontents{toc}{\protect\thispagestyle{empty}} % unterdrückt Seitenzahl
\addtocontents{lot}{\protect\thispagestyle{empty}} % unterdrückt Seitenzahl
\addtocontents{lof}{\protect\thispagestyle{empty}} % unterdrückt Seitenzahl

\tableofcontents			% Inhaltsverzeichnis



%% Der Text %%%%%%%%%%%%%%%%%%%%%%%%%%%%%%%%%%%%%%%%%%%%%%%%%%%%%%%%%%%
\frontmatter					% Vorspann (z.B. Glossar)

\phantomsection 

\printnoidxglossary[sort=letter, title=Abkürzungsverzeichnis, toctitle=Abkürzungsverzeichnis]

\mainmatter						% Hauptteil

\renewcommand{\baselinestretch}{1} %%% Zeilenabstand für das Dokument definieren.



\chapter{Einleitung}
\label{sec:Einleitung}

"`Augusta"' ist ein Chatbot, dessen Aufgabe es ist, dem Kunden Produkte aus dem Uni-Shop zu empfehlen. Dafür muss der Nutzer dem Chatbot via Texteingabe angeben, was gesucht wird. Diese Nutzereingaben werden mit Hilfe von Pattern-Matching analysiert und via PostgreSQL-Schnittstelle werden Antworten basierend auf Queries generiert, die dem geäußerten Wunsch des Nutzers entsprechen.\\
Das Ergebnis des Projekts kann auf GitHub unter \url{https://github.com/AlKYala/Chatbot} eingesehen werden.\\


\section{Ziele des Projekts}
\label{sec:ZielDesProjekts}

Vorrangiges Ziel des Projekts war der Versuch der Entwicklung eines computerlinguistischen Programms. Das Umsetzen von im Studium angeeigneter Fähigkeiten theoretischer und praktischer Natur, wie beispielsweise der Automatentheorie, Datenbankentechnologie und Softwaretechnik, hat eine essentielle Rolle in der Konzeption und Realisierung des Programms gespielt. Weiteres Ziel war das Einarbeiten in ein fremdes Framework mit eigenem Interpreter, hier ChatScript \citep{chatscript2019} von Bruce Wilcox, um ein entsprechendes Programm zu entwickeln.\\


\section{Der Chatbot Augusta}
\label{sec:Augusta}

Im Projektseminar der Computerlinguistik haben wir in kleinen Gruppen jeweils einen Chatbot für eine spezifische Anwendung entwickeln wollen. Unsere Gruppe, bestehend aus Ali Yalama und Julia Roegner, hat sich dazu entschieden, einen Chatbot als Kaufberatung für den Uni-Shop der Universität Trier zu bauen. Der Bot sollte für das Deutsche gelten und eine ausschließlich beratende Funktion haben, das bedeutet, dass der Bot Bestellungen oder Fragen zum Uni-Shop nicht beantworten können muss.\\
Als Basis dienten die Beispieldialoge (s. Anhang), die auch am Ende dieser Dokumentation abgedruckt sind. Hier wird davon ausgegangen, dass der Nutzer noch nicht weiß, welches Produkt er haben möchte, sondern ihm nur bekannt ist, wofür er es möchte, also entweder eine Geschenkidee hat oder etwas für eine bestimmte Anwendung sucht. Dann bietet der Chatbot dem Nutzer die Kategorien an, die in der Datenbank für Anwendungsgebiet oder Geschenkidee hinterlegt sind. Im Weiteren kann der Kunde dem Chatbot noch weitere Informationen über den gesuchten Artikel wie eine Farbe oder den Preis übergeben, am besten in Form eines kompakten Satzes, den der Chatbot analysiert. Dann nutzt der Bot die extrahierten Stichwörter, um in einer Datenbank, in der die Informationen über die Produkte des Uni-Shops gespeichert sind, nach etwas Passendem zu suchen.\\
Der gefundene Gegenstand wird dem Kunden angeboten und der Kunde kann entscheiden, ob er ihn ablehnt, dann sucht der Bot erneut, oder ihn annimmt, woraufhin der Chatbot fragt, ob eine neue Suche mit anderen Kriterien durchführen soll. Wenn nicht, beendet sich der Bot.\\
Der Name des Bots, Augusta, leitet sich vom Namen der antiken Stadt \textit{Augusta Treveorum} an der Mosel her, aus der das heutige Trier hervorgegangen ist.\\


\section{Struktur der Dokumentation}
\label{sec:Inhaltsbeschreibung} 
In Kapitel \ref{sec:InstallationAugusta} werden die einzelnen Versionen Augustas erläutert sowie die Installation beschrieben. In Kapitel \ref{sec:ChatScript: Grundlagen} folgen dann einige grundlegende Regeln zum Schreiben von Chatbots in ChatScript, woran sich ein Kapitel zu der Programmlogik hinter unserem Chatbot anschließt, das Augusta als deterministischen Automaten vorstellt. Informationen zu der Datenbank \textit{Uni-Shop}, in der die Produktinformationen der Artikel gespeichert sind, und ihrer Einbindung sowie der Suchanfragen finden sich in Kapitel \ref{sec:DBuEinbindung}. Auf die Benutzerschnittstelle wird in Kapitel \ref{sec:Benutzerschnittstelle} eingegangen. In Kapitel \ref{sec:Probleme} werden schließlich einige Probleme erläutert, auf die wir gestoßen sind und im Rahmen des Projekts leider nicht mehr lösen konnten.\\
Am Ende der Dokumentation sind die Dialoge abgedruckt, die als Basis dienten, und die Zusammenfassung des Ergebnisses aus einem kleinen Usability Test.

\chapter{ChatScript: Grundlagen}
\label{sec:ChatScript: Grundlagen}

Bei ChatScript handelt es sich um ein Framework, der es Entwicklern erlaubt, regelbasierte Chatbots zu entwickeln. Es handelt sich hierbei um "Natural Language tool" mit eigenem Interpreter, welcher Code in C++ übersezt. 


 

\section{ChatScript: Regeln}
\label{sec:ChatScript: Regeln}

\section{ChatScript: Topics}
\label{sec:ChatScript: Topics}

Die Quelldateien in einem ChatScript-Programm werden auch als "Topics" bezeichnet und haben die Dateiendung .top. Topics können als Zusammenfassung mehrerer Regeln betrachtet werden. Es bietet sich an, Topics als Module eines Chatbots zu handhaben. So wird z.B. in diesem Projekt "Augusta" die Datenbankabfrage als eignenes Topic gehandhabt. 

\section{ChatScript: Konzepte}
\label{sec:ChatScript: Konzepte}

\section{ChatScript: Bedingungen zur Steuerung von Programmablauf}
\label{sec:ChatScript: Bedingungen zur Steuerung von Programmablauf}

\section{ChatScript: Befehle zur Steuerung von Programmablauf}
\label{sec:ChatScript: Befehle zur Steuerung von Programmablauf}

\section{ChatScript: Pattern-Matching}
\label{sec:ChatScript: Pattern-Matching}

\section{ChatScript: Zufällige Ausgabe}
\label{sec:ChatScript: Zufällige Ausgabe}

\chapter{Augusta: Programmlogik}
\label{sec:Augusta: Programmlogik}



Hier, stelle in z.B. einem Automaten dar, wie der Programmablauf normalerweise ablaufen soll
\section{Befehle und Programmablauf}
\label{sec:Befehle und Programmablauf}

\section{Variablen und Programmablauf}
\label{sec:Variablen und Programmablauf}

\section{Konzepte und Programmablauf}
\label{sec:Konzepte und Programmablauf}

\section{Programmablauf im Detail}
\label{sec:Programmablauf im Detail}


\chapter{Kapitel1j}
\label{sec:Kapitel1j}

\section{UnterKapitel2j}
\label{sec:UnterKapitel2j}

\section{UnterKapitel3j}
\label{sec:UnterKapitel3j}
\chapter{Fazit}
\label{sec:Fazit}


%Glossar ausgeben
%\printglossary[style=altlist,title=Glossar]
\printglossary

\listoftables				% Tabellenverzeichnis
\listoffigures			% Abbildungsverzeichnis

\newpage

%\setbibpreamble{Präambel}		% Text vor dem Verzeichnis
%\bibliographystyle{dinat}
%\bibliography{bibliographie}	% Sie benötigen einen *.bib-Datei

%%%%%%%%%%%%%%%%%%%%%%%%%%%%%%%%%%%%%%%%%%%%%%%%
% B I B L I O G R A P H I E
%%%%%%%%%%%%%%%%%%%%%%%%%%%%%%%%%%%%%%%%%%%%%%%%

\renewcommand{\baselinestretch}{1}\normalsize 	%Zeilenabstand innerhalb der Angaben
\cleardoublepage
\addcontentsline{toc}{chapter}{\bibname}	%Bib. im Inhaltsvz.

\bibliography{CL-bib}					%Bibliographiedatenbank
\bibliographystyle{CL-DE}		%Zitierstruktur - natbib muss deklariert sein!

\nocite{*} 		%optional - gibt die gesammte Bib-Datenbank aus ansonsten werden nur 
%die zitierten Werke ausgegeben

%%%%%%%%%%%%%%%%%%%%%%%%%%%%%%%%%%%%%%%%%%%%%%%%
\end{document}
